% Options for packages loaded elsewhere
\PassOptionsToPackage{unicode}{hyperref}
\PassOptionsToPackage{hyphens}{url}
\PassOptionsToPackage{space}{xeCJK}
%
\documentclass[
]{article}
\usepackage{amsmath,amssymb}
% \usepackage{lmodern}
\usepackage{iftex}
\ifPDFTeX
  \usepackage[T1]{fontenc}
  \usepackage[utf8]{inputenc}
  \usepackage{textcomp} % provide euro and other symbols
\else % if luatex or xetex
  \usepackage{unicode-math}
  \defaultfontfeatures{Scale=MatchLowercase}
  \defaultfontfeatures[\rmfamily]{Ligatures=TeX,Scale=1}
  \setmainfont[]{DejaVu Serif}
  \setmathfont[]{Asana Math}
  \ifXeTeX
    \usepackage{xeCJK}
    \setCJKmainfont[]{Noto Serif CJK TC}
  \fi
  \ifLuaTeX
    \usepackage[]{luatexja-fontspec}
    \setmainjfont[]{Noto Serif CJK TC}
  \fi
\fi
% Use upquote if available, for straight quotes in verbatim environments
\IfFileExists{upquote.sty}{\usepackage{upquote}}{}
\IfFileExists{microtype.sty}{% use microtype if available
  \usepackage[]{microtype}
  \UseMicrotypeSet[protrusion]{basicmath} % disable protrusion for tt fonts
}{}
\makeatletter
\@ifundefined{KOMAClassName}{% if non-KOMA class
  \IfFileExists{parskip.sty}{%
    \usepackage{parskip}
  }{% else
    \setlength{\parindent}{0pt}
    \setlength{\parskip}{6pt plus 2pt minus 1pt}}
}{% if KOMA class
  \KOMAoptions{parskip=half}}
\makeatother
\usepackage{xcolor}
\IfFileExists{xurl.sty}{\usepackage{xurl}}{} % add URL line breaks if available
\IfFileExists{bookmark.sty}{\usepackage{bookmark}}{\usepackage{hyperref}}
\hypersetup{
  hidelinks,
  pdfcreator={LaTeX via pandoc}}
\urlstyle{same} % disable monospaced font for URLs
\setlength{\emergencystretch}{3em} % prevent overfull lines
\providecommand{\tightlist}{%
  \setlength{\itemsep}{0pt}\setlength{\parskip}{0pt}}
\setcounter{secnumdepth}{-\maxdimen} % remove section numbering
\usepackage[margin=1in]{geometry}
\usepackage{microtype}
\usepackage{hyperref}
\hypersetup{colorlinks=true, linkcolor=blue, urlcolor=blue}

\usepackage{fontspec}
\usepackage{xeCJK}

\setmainfont{DejaVu Serif}
\setCJKmainfont{Noto Serif CJK TC}[BoldFont = Noto Serif CJK TC]

\usepackage{amsmath, amssymb, amsthm}
\ifLuaTeX
  \usepackage{selnolig}  % disable illegal ligatures
\fi

\author{}
\date{}

\begin{document}

在 HORS 裡，我們會將每個區塊所對應到的私鑰 sk\_index
傳送給驗章者，驗章者再來檢查是否有 \[
\mathcal{H}(\mathrm{sk}_{\mathrm{index}} )= \mathrm{pk}_{\mathrm{index}}
\]

這樣的話，假設每個區段長度為 log2(t) 位元，那驗章者需要儲存 t
個公鑰。而其中 t 是二的某次方，並不小。這讓我們思考，是否可以使用 Merkle
tree 的技術，讓驗章者只需儲存少量公鑰即可？

今天介紹兩個把 Merkle tree 跟 HORS 結合的簽章系統。

\hypertarget{horst-ux7c3dux7ae0ux7cfbux7d71}{%
\section{HORST 簽章系統}\label{horst-ux7c3dux7ae0ux7cfbux7d71}}

\hypertarget{ux53c3ux6578ux8207ux9470ux5319ux751fux6210}{%
\subsection{參數與鑰匙生成}\label{ux53c3ux6578ux8207ux9470ux5319ux751fux6210}}

\hypertarget{ux53c3ux6578}{%
\subsubsection{參數}\label{ux53c3ux6578}}

參數 t ：我們會把訊息分割為 log2(t) 的長度 \[
d = \underbrace{ d_{1} }_{ \log_{2}t },\underbrace{ d_{2} }_{ \log_{2}t},\dots,\underbrace{ d_{end} }_{ \log_{2}t}
\] 我們設定所使用的雜湊函數是使用 SHA-256
，也設定這個訊息是雜湊函數後的結果（因此固定長度）

\hypertarget{ux516cux79c1ux9470}{%
\subsubsection{公私鑰}\label{ux516cux79c1ux9470}}

私鑰 \[
\mathrm{SK}=\{\mathrm{sk}_{0},\mathrm{sk}_{1},\dots,\mathrm{sk}_{t-1}\}
\] 其中每個 sk 都是長度為 n = 256 的二進位字串。 \[
\mathrm{PK}=\{\mathrm{pk}_{0},\mathrm{pk}_{1},\dots,\mathrm{pk}_{t-1}\}
\] \[
\mathrm{pk}_{i}= \mathcal{H}(\mathrm{sk}_{i})
\]

因為我們選用 SHA-256 所以每個 pk 都是 n = 256 長度的二進位字串。

此時，我們把這一堆公鑰 pk\_0, \ldots{} pk\_(t-1) 放入一個 Merkle tree
的最底層，並照 Merkle tree 的規則來算出 root 。將這個 root 發送為公鑰。
\#\# 簽章 首先先進行訊息分段： \[
d = \underbrace{ d_{1} }_{ \log_{2}t },\underbrace{ d_{2} }_{ \log_{2}t},\dots,\underbrace{ d_{end} }_{ \log_{2}t}
\] 計算每個區塊的索引 \[
\mathrm{index}_{i} = \mathrm{int}(d_{i},2)
\] \[
\mathrm{indices} = \{ \mathrm{index}_{1},\mathrm{index}_{2},\dots, \mathrm{index}_{end}\}
\]

並找出每段對應的私鑰 sk\_index\_i

除此之外，我們也提供從每個節點 pk\_index\_i 到 Merkle root 的必要材料
Auth\_i 於是

\[
\sigma_{i} = (\mathrm{sk}_{\mathrm{index}_{i}} ,\mathrm{pk}_{\mathrm{index}_{i}}, Auth_{i})
\] 發送完整的簽名 \[
\sigma 
= \{\sigma_{1},\dots,\sigma_{end}\}
\]

\hypertarget{ux9a57ux7ae0}{%
\subsection{驗章}\label{ux9a57ux7ae0}}

首先先進行訊息分段： \[
d = \underbrace{ d_{1} }_{ \log_{2}t },\underbrace{ d_{2} }_{ \log_{2}t},\dots,\underbrace{ d_{end} }_{ \log_{2}t}
\] 計算每個區塊的索引 \[
\mathrm{index}_{i} = \mathrm{int}(d_{i},2)
\] \[
\mathrm{indices} = \{ \mathrm{index}_{1},\mathrm{index}_{2},\dots, \mathrm{index}_{end}\}
\] 驗證所有索引的 hash 值 \[
\mathcal{H}(\sigma_{i}) 
\overset{?}{=} \mathrm{pk}_{\mathrm{index}_{i}} 
\] 除此之外，我們也需驗證 pk\_index\_i 是否屬於總公鑰 pk 的 Merkle tree
裡。若兩個條件都達成，則此簽章為有效。

\hypertarget{fors-ux7c3dux7ae0ux7cfbux7d71}{%
\section{FORS 簽章系統}\label{fors-ux7c3dux7ae0ux7cfbux7d71}}

\hypertarget{ux53c3ux6578ux8207ux9470ux5319ux751fux6210-1}{%
\subsection{參數與鑰匙生成}\label{ux53c3ux6578ux8207ux9470ux5319ux751fux6210-1}}

\hypertarget{ux53c3ux6578-1}{%
\subsubsection{參數}\label{ux53c3ux6578-1}}

參數 t ：我們會把訊息分割為 log2(t) 的長度 \[
d = \underbrace{ d_{1} }_{ \log_{2}t },\underbrace{ d_{2} }_{ \log_{2}t},\dots,\underbrace{ d_{end} }_{ \log_{2}t}
\] 我們設定所使用的雜湊函數是使用 SHA-256
，也設定這個訊息是雜湊函數後的結果（因此固定長度）。假設訊息被分段為 k
段。

\hypertarget{ux516cux79c1ux9470-1}{%
\subsubsection{公私鑰}\label{ux516cux79c1ux9470-1}}

我們會對每一個區塊都生成 w 個私鑰 \[
\mathrm{SK}_{1}=\{\mathrm{sk}_{1,0},\mathrm{sk}_{1,1},\dots,\mathrm{sk}_{1,t-1}\}
\] \[
\mathrm{SK}_{2}=\{\mathrm{sk}_{2,0},\mathrm{sk}_{2,1},\dots,\mathrm{sk}_{2,t-1}\}
\] \[
\dots
\] \[
\mathrm{SK}_{k}=\{\mathrm{sk}_{k,0},\mathrm{sk}_{k,1},\dots,\mathrm{sk}_{k,t-1}\}
\] 其中每個 sk\_\{i,j\} 都是長度為 n = 256 的二進位字串。

對每一個 SK\_i 都形成他的公鑰 PK\_i \[
\mathrm{PK}_{i}=\{\mathrm{pk}_{i,0},\mathrm{pk}_{i,1},\dots,\mathrm{pk}_{i,t-1}\}
\] \[
\mathrm{pk}_{i,j}= \mathcal{H}(\mathrm{sk}_{i,j})
\] 然後對每一組 PK\_i 都形成他的 Merkle tree，於是一共會有 k 把Merkle
root 總公鑰 mpk\_1,\ldots,mpk\_k

\hypertarget{ux7c3dux7ae0}{%
\subsection{簽章}\label{ux7c3dux7ae0}}

首先先進行訊息分段： \[
d = \underbrace{ d_{1} }_{ \log_{2}t },\underbrace{ d_{2} }_{ \log_{2}t},\dots,\underbrace{ d_{end} }_{ \log_{2}t}
\] 計算每個區塊的索引 \[
\mathrm{index}_{i} = \mathrm{int}(d_{i},2)
\] \[
\mathrm{indices} = \{ \mathrm{index}_{1},\mathrm{index}_{2},\dots, \mathrm{index}_{end}\}
\]

並找出每段對應的私鑰，第 d\_i 段對應到 \[
\mathrm{sk}_{i,\mathrm{index}_{i}}
\]

除此之外，我們也提供從每個節點 pk\_index\_i 到 Merkle root 的必要材料
Auth\_i 於是

\[
\sigma_{i} = (\mathrm{sk}_{i,\mathrm{index}_{i}},\mathrm{pk}_{i,\mathrm{index}_{i}}, Auth_{i})
\]

發送完整的簽名 \[
\sigma 
= \{\sigma_{1},\dots,\sigma_{end}\}
\]

\hypertarget{ux9a57ux7ae0-1}{%
\subsection{驗章}\label{ux9a57ux7ae0-1}}

首先也先進行訊息分段： \[
d = \underbrace{ d_{1} }_{ \log_{2}t },\underbrace{ d_{2} }_{ \log_{2}t},\dots,\underbrace{ d_{end} }_{ \log_{2}t}
\] 計算每個區塊的索引 \[
\mathrm{index}_{i} = \mathrm{int}(d_{i},2)
\] \[
\mathrm{indices} = \{ \mathrm{index}_{1},\mathrm{index}_{2},\dots, \mathrm{index}_{end}\}
\] 驗證所有索引的 hash 值 \[
\mathcal{H}(\sigma_{i}) 
\overset{?}{=} \mathrm{pk}_{i,\mathrm{index}_{i}} 
\] 除此之外，我們也需驗證 pk\_\{i, index\_i\} 是否屬於總公鑰 mpk\_i 的
Merkle tree 裡。若兩個條件都達成，則此簽章為有效。

\end{document}
